\par\indent\SourceLeaf{202}%
Diese Integrationen lassen sich (die allgemeinen Schwierigkeiten des Integrirens bei Seite gesetzt) offenbar vor der Auflösung unserer Hauptaufgabe ausführen.

Wenn nun für \(T,U\) solche Functionen von \(t,u\) substituirt werden, wobei die Bedingung unsrer Hauptaufgabe erfüllt wird, so geht \(\Omega\) in \(mm\omega\) über, und es wird
\[
\frac{(\mathrm{d}P+i\,\mathrm{d}Q)\mathord{.}(\mathrm{d}P-i\,\mathrm{d}Q)}{(\mathrm{d}p+i\,\mathrm{d}q)\mathord{.}(\mathrm{d}p-i\,\mathrm{d}q)}=\frac{mmn}{N}
\]
\noindent Man sieht aber leicht, dass der Zähler im ersten Theile dieser Gleichung durch den Nenner nur dann theilbar sein kann, wenn
\[
\begin{aligned}
&\text{entweder }\mathrm{d}P+i\,\mathrm{d}Q\text{ durch }\mathrm{d}p+i\,\mathrm{d}q,\text{ und }\mathrm{d}P-i\,\mathrm{d}Q\text{ durch }\mathrm{d}p-i\,\mathrm{d}q,\\
&\text{oder }\mathrm{d}P+i\,\mathrm{d}Q\text{ durch }\mathrm{d}p-i\,\mathrm{d}q,\text{ und }\mathrm{d}P-i\,\mathrm{d}Q\text{ durch }\mathrm{d}p+i\,\mathrm{d}q
\end{aligned}
\]
\noindent theilbar ist. Im ersteren Falle wird demnach \(\mathrm{d}P+i\,\mathrm{d}Q\) verschwinden, wenn \(\mathrm{d}p+i\,\mathrm{d}q=0\), oder \(P+iQ\) wird constant werden, wenn \(p+iq\) constant angenommen wird, d. i. \(P+iQ\) wird bloss Function von \(p+iq\) sein, und eben so \(P-iQ\) Function von \(p-iq\). Im andern Falle wird \(P+iQ\) Function von \(p-iq\), und \(P-iQ\) Function von \(p+iq\) sein. Es ist leicht einzusehen, dass diese Folgerungen auch umgekehrt gelten, nemlich dass, wenn für \(P+iQ,P-iQ\) Functionen von \(p+iq,p-iq\) (entweder respective, oder verkehrt) angenommen werden. die endliche Theilbarkeit des \(\Omega\) durch \(\omega\), und sonach die oben erforderlich gefundene Proportionalität Statt haben wird.

Man überzeugt sich übrigens leicht, dass wenn z. B.
\[
\begin{aligned}
P+iQ&=f(p+iq)\\
P-iQ&=f'(p-iq)
\end{aligned}
\]
\noindent gesetzt werden, die Beschaffenheit der Function \(f'\) schon durch die von \(f\) bedingt wird. Wenn nemlich unter den constanten Grössen, welche letztere etwa involviren mag, keine andere als reelle befindlich sind, so wird die andere \(f'\) mit der \(f\) ganz identisch sein müssen, damit jedesmal reellen Werthen von \(p,q\) reelle Werthe von \(P,Q\) entsprechen; im entgegengesetzten Falle wird sich \(f'\) von \(f\) nur dadurch unterscheiden, dass in den imaginären Elementen von \(f\) statt \(i\) überall das entgegengesetzte \(-i\) gesetzt werden muss.

Man hat hiernächst
